\title{%
  Science in Security,
  wrapping up
}
\author{Daniel Bosk\thanks{%
    This material is available under the Creative Commons 
    Attribution-NonCommercial-ShareAlike (CC-BY-NC-SA) 4.0 international 
    license.
    The material was written with some aid from GitHub Copilot.
}}
\institute{%
  KTH EECS
}

\begin{frame}
  \maketitle
\end{frame}

\mode*

\begin{abstract}
  \input{abstract.tex}
\end{abstract}

\clearpage

\section{Security as a scientific pursuit}

\Textcite{SoKScienceOfSecurity} discusses security as a scientific pursuit.

\begin{frame}
  \begin{exercise}[Time to vent]
    Air your thoughts about the paper 
    \citetitle{SoKScienceOfSecurity}~\autocite{SoKScienceOfSecurity}.
  \end{exercise}

  \begin{center}
    \(<10\) minutes
  \end{center}
\end{frame}

\begin{frame}
  \begin{question}
    How did they\autocite{SoKScienceOfSecurity} know that this is so?
    What can we actually say from these results?
  \end{question}

  \begin{question}
    The paper\autocite{SoKScienceOfSecurity} is from 
    \citeyear{SoKScienceOfSecurity}.
    How well can we repeat this study to see what the state is today?
  \end{question}

  \begin{center}
    \(<10\) minutes
  \end{center}
\end{frame}

\begin{frame}
  \begin{question}
    What can empirical experiments contribute to security?
  \end{question}

  \begin{question}
    What can deductive reasoning contribute to security?
  \end{question}

  \begin{question}
    What can quantitative observations contribute to security?
  \end{question}

  \begin{question}
    What can qualitative observations contribute to security?
  \end{question}

  \begin{center}
    5 minutes
  \end{center}
\end{frame}

\begin{frame}
  \begin{remark}[Final take-away]
    Always look for \enquote{the full version of the paper}!
  \end{remark}
\end{frame}


\section{Some cases}

\begin{frame}[fragile]
  \begin{block}{The cases}
    \begin{itemize}
      \item Passwords
      \item Programming in not-C
      \item Adding features
      \item Mandatory Access Control (MAC)
      \item Multi-Factor Authentication (MFA)
      \item Firewalls
      \item Social Engineering Awareness Training
      \item Encrypt Data-at-Rest
      \item Regular Security Updates
      \item Penetration Testing (Pentesting)
    \end{itemize}
  \end{block}
\end{frame}

\begin{frame}[fragile]
  \begin{exercise}[Work out cases]
    \begin{itemize}
      \item Join the breakout room of your favourite case.
      \item Work out the \emph{methodology} for \emph{how to know}.
        \hfill 15 minutes% erase the space, otherwise these will be uneven
        \begin{itemize}
          \item Analyse and evaluate your different approaches.
          \item Try to settle on one.
        \end{itemize}
      \item Take a break.
        \hfill 10 minutes%
      \item Then you'll summarize in whole class.
        \hfill 50 minutes%
    \end{itemize}
  \end{exercise}

  \begin{question}[To answer in class]
    \begin{itemize}
      \item What did you conclude: how do you \emph{evaluate} if a solution 
        is correct?
      \item What were the disagreements?
    \end{itemize}
  \end{question}
\end{frame}

\begin{frame}
  \begin{block}{Passwords}
    To do password-based authentication, we're taught that we should salt
    and hash the passwords. The user submits the password to the server,
    the server adds a 128-bit salt (random number) and hashes the password
    using SHA256, then hashes the hash, and hashes that hash, and so on for 1000 
    times. That final hash is used to check if it's the same as the one stored. If 
    so, the user is accepted. Is this password-based authentication scheme secure?
  \end{block}

  \begin{question}
    \begin{itemize}
      \item What did you conclude: how do you \emph{evaluate} if this is 
        correct?
      \item What were the disagreements?
    \end{itemize}
  \end{question}
\end{frame}

\begin{frame}
  \begin{block}{Programming in not-C}
    When we write programs, we should avoid languages like C or C++. (This is even 
    suggested by some US government agencies.) It's better to use a language like 
    Rust, which has built-in memory safety. Writing programs in Rust will result in 
    fewer security vulnerabilities in the programs and will thus be more secure. Is 
    this really true?
  \end{block}

  \begin{question}
    \begin{itemize}
      \item What did you conclude: how do you \emph{evaluate} if this is 
        correct?
      \item What were the disagreements?
    \end{itemize}
  \end{question}
\end{frame}

\begin{frame}
  \begin{block}{Adding features}
    The Bank has used the same back-end software for 50 years (written in COBOL). 
    This software works and manages the Bank's accounts. Previously the Bank's 
    staff interacted with the software, e.g. to tranfer money between accounts.

    The Bank hired FutureSoftware to add a web-based front-end so that the 
    customers can interact with the software directly. That way customers can 
    transfer money between accounts without the need for Bank staff to do it for 
    them. (They wrote the software in Python.)

    The Bank then asked VETCYB consultants to evaluate if this is secure. How would 
    go do it?
  \end{block}

  \begin{question}
    \begin{itemize}
      \item What did you conclude: how do you \emph{evaluate} if this is 
        correct?
      \item What were the disagreements?
    \end{itemize}
  \end{question}
\end{frame}

\begin{frame}
  \begin{block}{Mandatory Access Control (MAC)}
    The implementation of Mandatory Access Control (MAC) mechanisms, such
    as the Bell-LaPadula Model, is considered crucial in environments that
    demand high levels of security and confidentiality. Organizations that
    handle sensitive or classified information must ensure that data access
    is strictly regulated. MAC models achieve this by enforcing control
    policies that limit access based on user clearance levels. Does this work?
  \end{block}

  \begin{question}
    \begin{itemize}
      \item What did you conclude: how do you \emph{evaluate} if this is 
        correct?
      \item What were the disagreements?
    \end{itemize}
  \end{question}
\end{frame}

\begin{frame}
  \begin{block}{Multi-Factor Authentication (MFA)}
    Multi-Factor Authentication (MFA) is promoted as an effective way to enhance 
    security. By requiring users to provide two or more verification factors before 
    granting access, MFA adds an additional layer of protection beyond just the 
    password. The methods of implementing MFA varies; such as SMS, email, or 
    authenticator apps. Is this secure?
  \end{block}

  \begin{question}
    \begin{itemize}
      \item What did you conclude: how do you \emph{evaluate} if this is 
        correct?
      \item What were the disagreements?
    \end{itemize}
  \end{question}
\end{frame}

\begin{frame}
  \begin{block}{Firewalls}
    Firewalls are a fundamental component of network security that are
    taught in introductory courses. They act as barriers between trusted and
    untrusted networks by controlling incoming and outgoing network traffic
    based on predetermined security rules. Firewalls can prevent
    unauthorized access and attacks. Is this true?
  \end{block}

  \begin{question}
    \begin{itemize}
      \item What did you conclude: how do you \emph{evaluate} if this is 
        correct?
      \item What were the disagreements?
    \end{itemize}
  \end{question}
\end{frame}

\begin{frame}
  \begin{block}{Social Engineering Awareness Training}
    Security courses often emphasize the importance of educating employees
    about social engineering attacks. The training aims to help employees
    recognize and avoid common tactics used in attacks, such as phishing.
    Is this secure? How effective, if at all, is such training in practice at 
    preventing social engineering breaches over time?
  \end{block}

  \begin{question}
    \begin{itemize}
      \item What did you conclude: how do you \emph{evaluate} if this is 
        correct?
      \item What were the disagreements?
    \end{itemize}
  \end{question}
\end{frame}

\begin{frame}
  \begin{block}{Encrypt Data-at-Rest}
    One common security practice is encrypting data-at-rest to protect
    sensitive information stored on devices or servers. By encrypting the
    data, unauthorized users are prevented from accessing it without the
    appropriate decryption key. Is this secure?
  \end{block}

  \begin{question}
    \begin{itemize}
      \item What did you conclude: how do you \emph{evaluate} if this is 
        correct?
      \item What were the disagreements?
    \end{itemize}
  \end{question}
\end{frame}

\begin{frame}
  \begin{block}{Regular Security Updates}
    It is standard advice to regularly update software to stay protected
    against newly discovered vulnerabilities. Updates often include patches
    for security flaws that could be exploited by attackers. Does this contribute 
    to increased security?
  \end{block}

  \begin{question}
    \begin{itemize}
      \item What did you conclude: how do you \emph{evaluate} if this is 
        correct?
      \item What were the disagreements?
    \end{itemize}
  \end{question}
\end{frame}

\begin{frame}
  \begin{block}{Penetration Testing (Pentesting)}
    Penetration testing, commonly referred to as pentesting, is a proactive
    approach to evaluating the security of a computer system, network,
    or application by simulating an attack from malicious outsiders and
    insiders. Pentesting can help identify vulnerabilities before attackers
    exploit them. A comprehensive pentest involves several stages, including
    planning, reconnaissance, scanning, exploitation, and reporting. After a 
    successful pentest the system will be secure. Is this true?
  \end{block}

  \begin{question}
    \begin{itemize}
      \item What did you conclude: how do you \emph{evaluate} if this is 
        correct?
      \item What were the disagreements?
    \end{itemize}
  \end{question}
\end{frame}
